\section{Limitations and planned Helmholtz integration}

The present evidence has the following explicit limitations.
\begin{itemize}
    \item The matrices are small, deterministic diagnostic examples. They do not
    represent the dimensions, spectra, nonnormality, or memory traffic of a
    production sparse wave problem.
    \item Only one compiler/platform run is reported. No cross-platform rounding
    study, sanitizer run, performance benchmark, or statistical ensemble is
    included.
    \item The solver is sequential and unpreconditioned. No claim is made about
    practical iteration counts for indefinite Helmholtz systems.
    \item Orthogonality is observed by internal test instrumentation and checked
    with extended-precision accumulation, but the experiment is not a formal
    rounding-error analysis.
    \item The true residual is recomputed in double precision. Severe cancellation,
    underflow, or overflow can still limit its reliability; detected nonfinite
    arithmetic is rejected, but all finite extreme cases cannot be covered.
    \item A small residual alone cannot certify a small solution error for an
    ill-conditioned system. Condition estimation is not implemented.
\end{itemize}

\subsection{Data required for iteration plots}

The current tests print the final residual estimate, the final independent wide
residual, the maximum estimate--solver-residual gap, and the maximum two-pass
orthogonality defect. The public result retains the solver's residual history
only in memory. The one-pass control returns a single final defect. These
aggregates do not determine iteration curves, so no such curves are inferred
from them in this manuscript.

An estimate-versus-true-residual plot would require exporting each observer
sample as a row containing the case, total iteration, restart cycle, local step,
$\widehat{\rho}_t$, $\rho_{\mathrm{solver},t}$, basis size, and breakdown flag.
A separate per-step wide residual would require access to each candidate
solution, which the present observer does not expose. The plot must label the
double-precision recomputation and the wide oracle distinctly.

An orthogonality plot would require collecting
$\lVert W_t^H W_t-I\rVert_F$ at every step of the production run and of the
local one-pass control, using the same Hilbert matrix, right-hand side, step
range, and basis-size convention. To isolate the effect of the MGS pass count,
the product and norm routines should also be matched. Raw CSV data and the
revision, compiler, flags, precision, and solver options should accompany any
figures. Zero values must be shown explicitly or identified as omitted on
logarithmic axes, not replaced by invented positive measurements. Producing
these traces requires additional test instrumentation and a repeated run; it
does not require changing the numerical solver algorithm.

\subsection{Integration with FDFD Helmholtz}

The next application target is the WaveCore five-point FDFD discretization of
$-\Delta u-k^2u=f$ with homogeneous Dirichlet boundaries. Integration should
first use manufactured solutions away from both continuous and discrete
resonances. Each run should record grid dimensions, $k$, restart length, tolerance,
iteration limit, solver status, complete residual history, independently
recomputed algebraic residual, and error against the manufactured solution.
Grid-refinement studies must separate discretization error from GMRES iteration
error. Near resonances, the experiment should report conditioning indicators and
must not infer field accuracy from residual size alone.

Only after this unpreconditioned baseline is verified should preconditioning be
introduced. Candidate preconditioners and performance changes will require
separate hypotheses, numerical-equivalence checks, and measurements; none are
evaluated in the present article.
